\documentclass{beamer}

\input{MathMacros}

\usepackage{beamerthemesplit}
\usepackage{enumerate}
\usepackage{amssymb}
\usepackage[all,cmtip]{xy}
\usepackage[mathscr]{eucal}
\usepackage[natbib=true,style=authoryear,backend=bibtex,useprefix=true]{biblatex}
\addbibresource{reading.bib}
\usepackage{graphicx,color}

\title{Security and Risk in LLM Systems: Data, Leakage, and Adversarial Use}
\author{Reuben Brasher}
\date{\today}

\begin{document}

\frame{\titlepage}

\section[Outline]{}
\frame{\tableofcontents}

\section{Lecture 6}

\frame{
\frametitle{LLMs as Data Systems}
LLMs must be understood as data-processing systems with security implications.
\begin{itemize}
\item They ingest and transform large volumes of data.
\item They may reproduce patterns from training data.
\item They blur boundaries between input and output.
\item They function as probabilistic data interfaces.
\item They introduce risks through normal usage.
\item They must be treated as data systems.
\end{itemize}
}

\frame{
\frametitle{From Failure to Security Risk}
Failures in LLMs naturally become security concerns.
\begin{itemize}
\item Incorrect outputs may expose sensitive information.
\item Hallucinations may misrepresent internal data.
\item Misuse may cause unintended disclosure.
\item Errors scale across many users.
\item Risk accumulates over time.
\item Reliability becomes governance.
\end{itemize}
}

\frame{
\frametitle{The Data Problem}
LLMs are trained on large datasets with limited transparency.
\begin{itemize}
\item Training sources are not fully disclosed.
\item Data may include sensitive material.
\item Provenance of information is unclear.
\item Users cannot inspect training data.
\item Boundaries of knowledge are unknown.
\item Exposure risk exists without direct input.
\end{itemize}
}

\frame{
\frametitle{Unknown Exposure Risk}
The model may reproduce data without visibility.
\begin{itemize}
\item Memorized fragments may appear.
\item Sensitive patterns may exist in weights.
\item Outputs may reflect training data.
\item Reproduction occurs under certain prompts.
\item Risk is probabilistic.
\item Verification is not possible.
\end{itemize}
}

\frame{
\frametitle{Data Leakage Through Usage}
User interaction creates additional exposure pathways.
\begin{itemize}
\item Conversations may be logged.
\item Inputs may be retained.
\item Users may share sensitive information.
\item Internal processes may be exposed.
\item Proprietary data may leave boundaries.
\item Data lifecycle is unclear.
\end{itemize}
}

\frame{
\frametitle{Input as Future Risk}
User inputs may influence future system behavior.
\begin{itemize}
\item Inputs may persist across sessions.
\item Data may enter training pipelines.
\item Patterns may be generalized.
\item Future outputs may reflect prior inputs.
\item Boundaries between input and training blur.
\item Users contribute to system knowledge.
\end{itemize}
}

\frame{
\frametitle{Model Inversion Attacks}
Attackers can extract data through systematic probing.
\begin{itemize}
\item Repeated queries reveal patterns.
\item Sensitive data can be reconstructed.
\item Outputs are refined iteratively.
\item Memorization can be exploited.
\item Statistical methods guide extraction.
\item The model becomes a target.
\end{itemize}
}

\frame{
\frametitle{Extraction Techniques}
Specific techniques increase attack effectiveness.
\begin{itemize}
\item Contextual escalation narrows targets.
\item Query saturation expands coverage.
\item Low-temperature sampling stabilizes output.
\item Prompt chaining refines results.
\item Repetition amplifies signals.
\item Iterative probing reconstructs data.
\end{itemize}
}

\frame{
\frametitle{Model as Probeable System}
The model behaves like a weakly protected database.
\begin{itemize}
\item It responds to structured queries.
\item It lacks strict access control.
\item It cannot detect user intent.
\item Information may leak across sessions.
\item Outputs expose latent knowledge.
\item Reconstruction occurs over time.
\end{itemize}
}

\frame{
\frametitle{Adversarial Use of LLMs}
LLMs enable malicious applications at scale.
\begin{itemize}
\item Phishing and scam messages can be generated.
\item Social engineering scripts can be created.
\item Malware development may be assisted.
\item Reconnaissance tasks can be automated.
\item Attack content scales easily.
\item Barriers to misuse are reduced.
\end{itemize}
}

\frame{
\frametitle{Dual-Use Capability}
LLMs support both beneficial and harmful uses.
\begin{itemize}
\item Security teams use models for analysis.
\item Attackers use models for automation.
\item Knowledge becomes widely accessible.
\item Expertise is commoditized.
\item Capabilities scale across users.
\item Risk increases with accessibility.
\end{itemize}
}

\frame{
\frametitle{Organizational Security Risk}
Organizations face direct exposure through usage.
\begin{itemize}
\item Confidential data may be leaked.
\item Internal workflows may be exposed.
\item Credentials may be revealed.
\item Data boundaries are difficult to enforce.
\item Attack surface increases.
\item Monitoring is incomplete.
\end{itemize}
}

\frame{
\frametitle{Compliance Risk}
LLMs introduce regulatory challenges.
\begin{itemize}
\item GDPR and CCPA violations may occur.
\item Data retention policies are unclear.
\item Data processing control is limited.
\item Auditing is difficult.
\item Unauthorized transfer may occur.
\item Legal liability increases.
\end{itemize}
}

\frame{
\frametitle{Trust Risk}
Model behavior affects organizational trust.
\begin{itemize}
\item Outputs may be incorrect.
\item Behavior may be inconsistent.
\item Decisions are hard to explain.
\item Users may overtrust outputs.
\item Failures damage reputation.
\item Confidence may degrade.
\end{itemize}
}

\frame{
\frametitle{Enterprise Requirements}
Organizations require strong governance properties.
\begin{itemize}
\item Systems must be auditable.
\item Data flows must be controlled.
\item Decisions must be traceable.
\item Behavior must be reproducible.
\item Accountability must be enforced.
\item Risk must be measurable.
\end{itemize}
}

\frame{
\frametitle{Fundamental Misalignment}
LLMs conflict with enterprise expectations.
\begin{itemize}
\item Internal processes are opaque.
\item Outputs are probabilistic.
\item No inherent audit trail exists.
\item Reproducibility is not guaranteed.
\item Control mechanisms are limited.
\item Explainability is weak.
\end{itemize}
}

\frame{
\frametitle{Policy Mitigation}
Policies reduce exposure risk.
\begin{itemize}
\item Restrict sensitive data entry.
\item Define confidential information categories.
\item Establish usage guidelines.
\item Enforce system boundaries.
\item Align with risk tolerance.
\item Require compliance standards.
\end{itemize}
}

\frame{
\frametitle{Training Mitigation}
User education is essential.
\begin{itemize}
\item Train users on risks.
\item Demonstrate failure cases.
\item Encourage cautious input.
\item Promote output verification.
\item Build awareness of exposure.
\item Develop literacy.
\end{itemize}
}

\frame{
\frametitle{Technical Controls}
Systems must enforce constraints programmatically.
\begin{itemize}
\item Disable data sharing when possible.
\item Monitor interactions.
\item Implement access controls.
\item Filter sensitive data.
\item Use secure environments.
\item Limit external exposure.
\end{itemize}
}

\frame{
\frametitle{Structural Insight}
Risks arise from system design.
\begin{itemize}
\item Data exposure is inherent.
\item Interaction creates risk.
\item Adversaries exploit properties.
\item Risks scale with usage.
\item Problems persist across contexts.
\item Mitigation must be systemic.
\end{itemize}
}

\frame{
\frametitle{Synthesis}
LLM risk is structural rather than incidental.
\begin{itemize}
\item It originates in training opacity.
\item It emerges through interaction.
\item It is amplified by adversaries.
\item It cannot be eliminated.
\item It must be managed continuously.
\item It requires layered controls.
\end{itemize}
}

\frame{
\frametitle{LLM Governance}
LLMs must be treated as governed systems.
\begin{itemize}
\item They process sensitive data.
\item They introduce measurable risk.
\item They require explicit controls.
\item They must meet enterprise standards.
\item They need monitoring.
\item They require oversight.
\end{itemize}
}

\frame{
\frametitle{Discussion}
LLMs must be integrated into enterprise governance frameworks.
\begin{itemize}
\item Policy frameworks are required.
\item Architectural constraints must exist.
\item Auditability is necessary.
\item Compliance integration is required.
\item Operational oversight is essential.
\item Systems become infrastructure.
\end{itemize}
}

\begin{frame}[t,allowframebreaks]
\frametitle{References}
\printbibliography
\end{frame}

\end{document}